\chapter{Gravedad: holonomía, torsión y respuesta colectiva}
\label{chap:gravedad}
\index{gravedad}\index{Cartan--Holst}\index{holonomía}
\index{torsión}\index{gravitón!no primitivo}

La gravedad no se introduce aquí mediante una fórmula decimal para \(G\).
Su construcción reúne dos ramas ya tipadas. En la primera, una ruta conserva
orientación y memoria; una conexión compara marcos en puntos distintos; su
holonomía mide el retorno y su curvatura, el defecto de cierre. En la segunda,
la carta dimensional construye \(G_{\rm int}(\theta_G)\) y
\(\kappa_{\rm int}\) a partir de la escala interna de acción, la velocidad y
la duración elemental. Sólo después de disponer de ambas ramas se escribe la
acción Palatini--Cartan--Holst. Ninguna cifra metrológica externa participa en
esa precedencia.

\section{Retornos: posición, orientación, fase y memoria}

Sea \(F_{\rm obs}\) la cronología observable del transporte y
\[
 \mathcal K_{27}=F_{\rm obs}^{27}.
\]
La jerarquía exacta distingue tres retornos:
\[
 \begin{array}{c|c}
 27\ \text{pasos}&
 \text{retorno de posición e inversión de orientación},\\
 54\ \text{pasos}&
 \text{retorno de posición y orientación},\\
 108\ \text{pasos}&
 \text{retorno completo de la fase observable}.
 \end{array}
\]
El retorno observable no borra el registro acumulado.  Si \(g(e)\) cuenta
una inversión de orientación y \(s(e)\) un cambio efectivo de hoja, entonces
\[
 c(\gamma)=\sum_{e\subset\gamma}(g(e),s(e))
\]
satisface la ley cocíclica de concatenación.  En los bloques canónicos,
\[
 c_{27}=(1,18),\qquad c_{108}=(4,72).
\]
Por tanto, la proyección observable de \(\mathcal K_{27}\) tiene orden
cuatro, mientras la coordenada acumulativa sigue creciendo.  Volver a la
misma fase no significa volver a la misma historia.

\begin{theorem}[Cobertura acumulativa del retorno]
\label{thm:cobertura-retorno-grav-rev6}
Sobre \(\mathcal O_{108}\times\mathbb N^2\), la elevación
\[
 \widetilde F(x,m)=
 \bigl(F_{\rm obs}x,m+c(x\to F_{\rm obs}x)\bigr)
\]
satisface
\[
 \widetilde{\mathcal K}_{27}(x,m)
 =\bigl(\mathcal K_{27}x,m+(1,18)\bigr)
\]
y
\[
 \widetilde{\mathcal K}_{27}^{\,4}(x,m)
 =\bigl(x,m+(4,72)\bigr)\neq(x,m).
\]
\end{theorem}

\begin{proof}
La constancia de \(c_{27}\) sobre la órbita y la ley de concatenación dan
\[
 c_{108}=\sum_{j=0}^{3}c_{27}(\mathcal K_{27}^jx)
 =4(1,18).
\]
La primera coordenada retorna; la segunda conserva la memoria acumulada.
\end{proof}

Éste es el dato que una realización geométrica debe preservar.  No se
identifica la memoria discreta con la torsión: se construye una representación
que permita compararlas.

\section{Del grupoide de caminos al grupo de Cartan}

Sea \(\mathcal G_{\rm mem}\) el grupoide de rutas enriquecidas, con
composición por concatenación e inversa por reversión. Fijemos un elemento
\(R_4\in\operatorname{Spin}^+(1,3)\) de orden cuatro y un vector temporal
unitario \(u\in\mathbb R^{1,3}\) fijado por \(R_4\). Si
\[
c(\gamma)=\bigl(c_g(\gamma),c_s(\gamma)\bigr)\in\mathbb Z^2
\]
es el cociclo aditivo de inversión y cambio de hoja, definimos
\[
\rho_{\rm C}^{\rm disc}(\gamma)
=
\left(
R_4^{\,c_g(\gamma)},
\ell_0c_s(\gamma)u
\right)
\in\operatorname{Spin}^+(1,3)\ltimes\mathbb R^{1,3}.
\]

\begin{theorem}[Realización discreta del cociclo de retorno]
\label{thm:realizacion-discreta-cartan-rev6}
\(\rho_{\rm C}^{\rm disc}\) es un funtor del grupoide de rutas en el grupo de
Cartan. En particular, conserva identidad, concatenación e inversión, y
transporta el retorno \(27\to54\to108\) sin borrar la coordenada acumulativa.
\end{theorem}

\begin{proof}
La ley del producto semidirecto es
\[
(R,a)(S,b)=(RS,a+Rb).
\]
Como \(R_4u=u\) y \(c\) es aditivo,
\[
\begin{aligned}
\rho_{\rm C}^{\rm disc}(\gamma_2)
\rho_{\rm C}^{\rm disc}(\gamma_1)
&=
\left(
R_4^{c_g(\gamma_2)+c_g(\gamma_1)},
\ell_0\bigl(c_s(\gamma_2)+c_s(\gamma_1)\bigr)u
\right)\\
&=\rho_{\rm C}^{\rm disc}(\gamma_2\gamma_1).
\end{aligned}
\]
La ruta identidad tiene cociclo nulo y la reversión cambia su signo, por lo
que se preservan identidad e inversa. Las fórmulas de retorno siguen de
\(c_{27}=(1,18)\) y \(c_{108}=4c_{27}\).
\end{proof}

Este teorema construye el puente discreto. Una realización suave posterior
consiste en una representación
\[
 \rho_{\rm C}:\mathcal G_{\rm mem}
 \longrightarrow\operatorname{Spin}(1,3)\ltimes\mathbb R^{1,3}
\]
y en campos \(e^I,\omega^{IJ}\) cuya holonomía entrelace esa representación.
El diagrama que fija el tipo es
\[
 \operatorname{Hol}_{\mathcal A_{\ell_0}}
   \bigl(\mathfrak R_{\rm C}(\gamma)\bigr)
 =
 \rho_{\rm C}\bigl(\operatorname{Hol}_{\rm TPK}(\gamma)\bigr).
\]
Identidades, inversión, concatenación y subdivisión deben conservarse.  Esta
ecuación no rebautiza el acarreo como curvatura: exige que el transporte
discreto ya representado y el geométrico realicen la misma composición.

Fijada una escala positiva \(\ell_0\), la conexión conjunta es
\[
 \mathcal A_{\ell_0}
 =
 \begin{pmatrix}
 \omega&\ell_0^{-1}e\\
 0&0
 \end{pmatrix}.
\]

\begin{proposition}[Curvatura conjunta de Cartan]
\label{prop:curvatura-conjunta-cartan-rev6}
La curvatura de \(\mathcal A_{\ell_0}\) es
\[
 \mathcal F_{\ell_0}
 =\dd\mathcal A_{\ell_0}
  +\mathcal A_{\ell_0}\wedge\mathcal A_{\ell_0}
 =
 \begin{pmatrix}
 F_\omega&\ell_0^{-1}T_{\rm tor}\\
 0&0
 \end{pmatrix},
\]
donde
\[
 F_\omega^{IJ}
 =\dd\omega^{IJ}+\omega^I{}_K\wedge\omega^{KJ},
 \qquad
 T_{\rm tor}^{I}=D_\omega e^I.
\]
\end{proposition}

\begin{proof}
Se multiplican las matrices de formas usando el producto exterior.  El
bloque diagonal produce \(F_\omega\); el bloque de traslación produce
\(\dd e+\omega\wedge e=D_\omega e\).
\end{proof}

Curvatura y torsión son dos componentes de un mismo defecto de transporte,
pero cumplen funciones distintas.  La curvatura mide el retorno de marcos;
la torsión mide el cierre de la co-tétrada.

\section{Gauss--Stokes y capacidad areal}

En un complejo celular finito, para toda \(1\)-cocadena \(\Omega\) y toda
\(2\)-cadena \(z\),
\[
 \boxed{\quad
 \langle\delta\Omega,z\rangle
 =\langle\Omega,\partial z\rangle.
 \quad}
\]
La identidad es la definición del coborde como adjunto de la frontera.
Cuando \(z\) suma las caras coherentemente orientadas de una
pseudovariedad, las aristas interiores se cancelan y sólo permanece el
borde geométrico.  La misma igualdad se leyó ya como
vorticidad--circulación y como curvatura--holonomía.

La escala areal también posee un modelo exacto.  En el grafo cúbico
\(N\times N\times N\), todo corte entre dos caras opuestas contiene al menos
\(N^2\) aristas y una capa transversal alcanza la cota.  Para \(N=9\),
\[
 9^3=729=27^2.
\]
Esta igualdad cardinal relaciona volumen y corte mínimo; no identifica los
grupos \((\mathbb Z/9)^3\) y \((\mathbb Z/27)^2\), que tienen exponentes
distintos.

\section{La acción Palatini--Cartan--Holst}

Sea \(M\) una variedad lisa orientada de dimensión cuatro, con co-tétrada no
degenerada \(e^I\), conexión de Lorentz
\(\omega^{IJ}=-\omega^{JI}\), y condiciones de borde que anulen el término
variacional de frontera.  Escribamos
\[
 \Sigma^{IJ}=e^I\wedge e^J,\qquad
 (\star X)^{IJ}=\frac12\epsilon^{IJ}{}_{KL}X^{KL},
 \qquad
 G_{\rm int}=G_{\rm int}(\theta_G),\qquad
 \kappa_{\rm int}=\frac{8\pi G_{\rm int}}{c_{\rm int}^{4}}.
\]
Estas dos últimas cantidades no se definen mediante la acción: proceden de
la construcción constitutiva de la \cref{eq:G-interno}. El
\cref{thm:homogeneidad-pch} demuestra antes que la combinación siguiente
pertenece a la recta de acción.
La acción de primer orden es
\[
\boxed{
 S[e,\omega,\Psi]
 =\frac1{2\kappa_{\rm int}c_{\rm int}}\int_M
 \Sigma_{IJ}\wedge
 \left(\star F^{IJ}+\frac1\gamma F^{IJ}\right)
 -\frac{\Lambda_{\rm int}}{\kappa_{\rm int}c_{\rm int}}
  \int_M\operatorname{vol}_e
 +S_{\rm m}[e,\omega,\Psi].}
\]
El parámetro interno construido en el capítulo anterior especializa
\[
 \gamma=\Gamma_{6\to8}.
\]
La sustitución se efectúa después de obtener
\(\Gamma_{6\to8}\) desde \(A,C^\ast,q_\pm\) y la inclusión
hexada--octada; la acción no se usa para redefinir esas entradas.

\begin{theorem}[Ecuación de Cartan--Holst]
\label{thm:cartan-holst-rev6}
Si la corriente de espín se normaliza mediante
\[
 \delta_\omega S_{\rm m}
 =-\frac12\int_M\delta\omega^{IJ}\wedge\sigma_{IJ},
\]
la variación respecto de \(\omega\) produce
\[
 \boxed{\quad
 D_\omega
 \left[(\star+\gamma^{-1}I)\Sigma^{IJ}\right]
 =\kappa_{\rm int}c_{\rm int}\,\sigma^{IJ}.
 \quad}
\]
\end{theorem}

\begin{proof}
Se usa \(\delta F^{IJ}=D_\omega\delta\omega^{IJ}\), se integra por partes y
se iguala a cero el coeficiente de la variación arbitraria.  El término
cosmológico no depende de \(\omega\).
\end{proof}

En firma lorentziana, \(\star^2=-I\) sobre bivectores.  Para
\(\gamma\in\mathbb R\setminus\{0\}\),
\[
 (\star+\gamma^{-1}I)^{-1}
 =\frac{\gamma^2}{1+\gamma^2}
  (-\star+\gamma^{-1}I).
\]
Por ello, si la corriente de espín se anula,
\[
 T_{\rm tor}^I=0,\qquad
 \omega=\mathring\omega(e).
\]
Con materia espinorial, la torsión se determina algebraicamente por la
corriente.  No constituye, en la acción mínima, un grado propagante
independiente.

\section{Ecuación métrica, espín y balance geométrico}

Para fijar sin ambigüedad la normalización, definimos las fuentes medidas en
la recta de acción por
\[
 \delta_g S_{\rm m}
 =-\frac12\int_M\operatorname{vol}_e\,
 \mathcal T^{S}_{\mu\nu}\,\delta g^{\mu\nu}.
\]
Cuando las coordenadas de la tétrada tienen dimensión de longitud, escribimos
\[
 T^{\rm phys}_{\mu\nu}:=c_{\rm int}\mathcal T^{S}_{\mu\nu},
 \qquad
 U^{\rm phys,spin}_{\mu\nu}
 :=c_{\rm int}\mathcal U^{S,\rm spin}_{\mu\nu}.
\]
Tras eliminar algebraicamente la conexión, la ecuación métrica adopta las
dos formas equivalentes
\[
 \boxed{
 G_{\mu\nu}+\Lambda_{\rm int}g_{\mu\nu}
 =\kappa_{\rm int}c_{\rm int}
 \bigl(\mathcal T_{\mu\nu}^{S,\rm LC}
       +\mathcal U_{\mu\nu}^{S,\rm spin}\bigr)
 =\kappa_{\rm int}
 \bigl(T_{\mu\nu}^{\rm phys,LC}
       +U_{\mu\nu}^{\rm phys,spin}\bigr).}
\]
La contribución de espín es cuadrática en la corriente correspondiente y
depende del acoplamiento de materia; su coeficiente no se fija sin declarar
la acción fermiónica y las convenciones de dualidad.

Las identidades geométricas
\[
 D_\omega T_{\rm tor}^{I}=F^{I}{}_{J}\wedge e^{J},
 \qquad
 D_\omega F^{IJ}=0
\]
implican, después de eliminar la torsión,
\[
 \nabla^\mu\bigl(T_{\mu\nu}^{\rm phys,LC}
             +U_{\mu\nu}^{\rm phys,spin}\bigr)=0.
\]
Este balance de Bianchi--Noether es una condición de consistencia del
transductor gravitatorio: ningún término adicional puede incorporarse sin
proceder de una acción compatible o satisfacer el mismo balance.

\section{Correspondencia areal–volumétrica \(\pi/\varphi^2\)}

El calendario \((\Sigma,\Pi,\Pi)\) produce la matriz de incidencia
\[
 F_{\rm av}=\begin{pmatrix}0&1\\1&1\end{pmatrix}.
\]
Su medida de Parry da
\(\mu_{\rm ar}/\mu_{\rm vol}=\varphi^{-2}\), y el lector regional marca una
sola vez el cierre \(\pi_{\rm HMT}\).  El retorno marcado satisface
\[
 \Theta_n
 =\pi_{\rm HMT}\frac{F_{n-1}}{F_{n+1}}
 \longrightarrow
 \frac{\pi_{\rm HMT}}{\varphi^2}.
\]
Esta razón no es otra versión de \(G\). Es la correspondencia adimensional entre
lectura areal y lectura volumétrica.  Bajo una misma amplitud positiva
\(\mathcal M\),
\[
 \mathcal I_{\rm av}=\mu_{\rm vol}\mathcal M,\qquad
 \mathcal G_{\rm av}=\pi_{\rm HMT}\mu_{\rm ar}\mathcal M
\]
dan exactamente
\[
 \frac{\mathcal G_{\rm av}}{\mathcal I_{\rm av}}
 =\frac{\pi_{\rm HMT}}{\varphi^2}.
\]
En la hoja plana, la restricción diagonal es
\(\mathcal G_{\rm tor}=\mathcal I_{\rm tor}\).  La equivalencia local y la
ambivalencia de frontera son, por tanto, dos regímenes del lector, no dos
valores incompatibles de una masa.

\section{Tipado de \(G\) como transductor gravitatorio y orden causal de sus lectores}

El tipo y la homogeneidad de \(G\) pertenecen a la construcción
constitutiva previa. La \cref{eq:G-interno} fija una sola vez la familia
covariante \(G_{\rm int}(\theta_G)\), y el
\cref{thm:homogeneidad-pch} demuestra que su inserción en la acción posee
dimensión de acción. Este capítulo no redefine esa familia: utiliza su
sección junto con la holonomía ya construida. La causalidad no es lineal,
sino convergente:
\[
\begin{aligned}
 \mathcal B_{\rm geom}&:
 \text{retorno y cociclo}\to\text{representación de Cartan}
 \to(e,\omega),\\
 \mathcal B_{\rm dim}&:
 (\hbar_{\rm int},c_{\rm int},t_0,\theta_G)
 \to G_{\rm int}\to\kappa_{\rm int},
\end{aligned}
\]
\[
 (\mathcal B_{\rm geom},\mathcal B_{\rm dim})
 \to S_{\rm PCH}
 \to\text{torsión algebraica, ecuación métrica y balance}.
\]
La comparación metrológica pertenece al final y no selecciona ninguna de
las dos ramas.

\Needspace{22\baselineskip}
\section{Origen colectivo del modo gravitatorio en la holonomía de rutas}

La acción anterior tiene como objetos básicos la co-tétrada y la conexión.
La torsión de Cartan es algebraica; la curvatura posee el sector propagante que
contiene las ondas gravitatorias.  Por ello el formalismo no necesita
postular un gravitón elemental para definir la interacción.  Una eventual
cuantización de las perturbaciones puede producir un modo colectivo de
helicidades \(\pm2\); ese modo sería una excitación del transporte
geométrico, no un dato añadido al nivel de APP, TRIT o TPK.

Esta distinción también excluye una confusión frecuente.  La existencia de
ondas clásicas demuestra propagación de curvatura; no demuestra por sí sola
un cuanto fundamental.  Recíprocamente, describir el campo mediante
geometría no elimina sus grados radiativos.  HMT--MD sitúa la pregunta en el
lugar correcto: primero se construye la memoria y su holonomía; después se
estudia el espectro de sus excitaciones.
